\documentclass[12pt]{article}
\usepackage[a4paper,margin=1in]{geometry}
\usepackage[T1]{fontenc}
\usepackage{lmodern,microtype}
\usepackage{amsmath,amssymb,amsthm,mathtools}
\usepackage{booktabs,array,enumitem,xcolor}
\usepackage[numbers,sort&compress]{natbib}
\usepackage[colorlinks=true,linkcolor=blue!55!black,citecolor=blue!55!black,urlcolor=blue!55!black]{hyperref}
\linespread{1.07}

\newtheorem{theorem}{Theorem}[section]
\newtheorem{proposition}[theorem]{Proposition}
\newtheorem{corollary}[theorem]{Corollary}
\theoremstyle{remark}
\newtheorem{remark}[theorem]{Remark}
\newcommand{\cH}{\mathcal H}
\newcommand{\norm}[1]{\left\lVert#1\right\rVert}

\title{The Projective Wrap Obstruction\\
\large Exact Phase and Scalar Lower Bounds for Finite Temporal Cycles}
\author{Liang Wang}
\date{July 2026}

\begin{document}
\maketitle

\begin{abstract}
RH-182 found no orthogonal finite temporal clock passing a joint wrap,
primal, and adjoint gate.  A possible concern is that this negative result
could be an artifact of a sign, phase convention, or poorly normalized final
edge.  This paper removes that ambiguity.

For unit vectors $x_0,x_L$ in a complex Hilbert space, the best unimodular
wrap has phase
$\omega=\langle x_0,x_L\rangle/|\langle x_0,x_L\rangle|$ and chord
$\sqrt{2-2|\langle x_0,x_L\rangle|}$.  If the final edge is allowed an
arbitrary complex scalar, the exact minimum residual is the physical edge
amplitude times
$\sqrt{1-|\langle x_0,x_L\rangle|^2}$.  Hence the projective return distance
is an unavoidable lower bound for every one-edge cyclic closure with the
declared temporal span.  For a polar-orthogonalized synthesis, the residual
is bounded above and below by the same rank-one defect divided by the extreme
square roots of its Gram spectrum.

Eighty deterministic formula cases verify the rank-one, scalar-optimal, and
polar-bound identities with zero failures; the largest direct formula error
is $4.45\times10^{-16}$.  Replaying the 126 RH-182 physical windows, 43
negative orientation marks improve the endpoint chord, by as much as
$1.5849$.  Nevertheless no projective return distance is at most $0.25$ and
no orthogonal three-way gate survives.

This gives a rigorous finite obstruction for the declared temporal spans:
phase repair cannot rescue the orthogonal branch.  It is not an all-level
no-go theorem and does not constrain a construction with distinct right and
left spaces.  No physical Riesz or Hilbert--Polya conclusion is claimed.
\end{abstract}

\section{Why a separate wrap theorem is needed}

The weighted temporal clock of RH-182 has exact chain columns.  Its only
intertwining error is the endpoint-to-seed wrap \citep{WangRH182}.  In the
real physical data, many length-four endpoint correlations are negative.
A positive wrap edge therefore overstates the error.  The orientation mark
corrects that sign, but one must still know whether some more general phase
or scalar can eliminate the defect.

The answer is exact and elementary.  It is useful because it separates
three increasingly permissive questions:
\begin{enumerate}[leftmargin=2em]
\item positive wrap: compare $x_L$ with $x_0$;
\item phase-marked wrap: compare $x_L$ with $\omega x_0$,
$|\omega|=1$;
\item arbitrary scalar wrap: compare $x_L$ with $\beta x_0$,
$\beta\in\mathbb C$.
\end{enumerate}
If the third question still has a large residual, neither sign nor amplitude
normalization is the problem.

\section{Phase optimization}

Let $u,v\in\cH$ be unit vectors and put $c=\langle u,v\rangle$.

\begin{theorem}[Optimal phase wrap]\label{thm:phase}
If $c\ne0$, the minimizer of
$\norm{v-\omega u}$ over $\omega\in\mathbb T$ is
\begin{equation}\label{eq:optimal-phase}
 \omega_*=\frac{c}{|c|}.
\end{equation}
The minimum is
\begin{equation}\label{eq:phase-chord}
 d_{\mathbb T}(u,v)
 =\sqrt{2-2|c|}.
\end{equation}
If $c=0$, every phase is minimizing and the minimum is $\sqrt2$.
\end{theorem}

\begin{proof}
For $|\omega|=1$,
\[
 \norm{v-\omega u}^2
 =2-2\operatorname{Re}(\overline\omega c).
\]
The real part is at most $|c|$, with equality at
$\omega=c/|c|$.  Substitution gives \eqref{eq:phase-chord}.
\end{proof}

The quantity in \eqref{eq:phase-chord} is a chord distance on projective
space.  For small principal angle it is comparable to that angle, but the
exact formula is retained in the audit.

\section{Arbitrary scalar optimization}

Allowing an arbitrary scalar gives the orthogonal projection of $v$ onto
$\mathbb Cu$.

\begin{theorem}[Projective scalar obstruction]\label{thm:scalar}
For $a\ge0$,
\begin{equation}\label{eq:scalar-min}
 \min_{\beta\in\mathbb C}\norm{av-\beta u}
 =a\sqrt{1-|\langle u,v\rangle|^2}.
\end{equation}
The minimizing scalar is
$\beta_*=a\langle u,v\rangle$.
\end{theorem}

\begin{proof}
Decompose
$v=\langle u,v\rangle u+v_\perp$ with $v_\perp\perp u$.  Then
\[
 av-\beta u
 =(a\langle u,v\rangle-\beta)u+av_\perp.
\]
The two summands are orthogonal.  The first vanishes exactly at
$\beta_*$, while
$\norm{v_\perp}^2=1-|\langle u,v\rangle|^2$.
\end{proof}

Define the projective return distance
\begin{equation}\label{eq:projective-distance}
 d_{\mathbb P}(u,v)=
 \sqrt{1-|\langle u,v\rangle|^2}.
\end{equation}
It is no larger than the phase chord, but unlike the phase chord it permits
amplitude adjustment.  A large value of \eqref{eq:projective-distance} is
therefore a stronger obstruction.

\begin{proposition}[Exact relation between the two distances]
For unit vectors $u,v$,
\begin{equation}\label{eq:distance-relation}
 d_{\mathbb T}(u,v)^2
 =\frac{2}{1+|\langle u,v\rangle|}
 d_{\mathbb P}(u,v)^2.
\end{equation}
Consequently
\begin{equation}\label{eq:distance-comparison}
 d_{\mathbb P}(u,v)\le d_{\mathbb T}(u,v)
 \le\sqrt2\,d_{\mathbb P}(u,v).
\end{equation}
\end{proposition}

\begin{proof}
Put $s=|\langle u,v\rangle|$.  Then
$d_{\mathbb T}^2=2(1-s)$ and
$d_{\mathbb P}^2=(1-s)(1+s)$.  Their ratio is $2/(1+s)$,
which lies in $[1,2]$.
\end{proof}

Equation \eqref{eq:distance-relation} quantifies exactly what is gained by
allowing amplitude as well as phase.  The gain is never more than a factor
$\sqrt2$ in norm.  Hence a projective failure cannot be attributed to a
minor convention about endpoint normalization.

\section{Application to one-edge temporal closure}

Let $x_j$ be the normalized orbit from RH-182 and let
$J=[x_t,\ldots,x_{t+L-1}]$.  The physical final edge is
\begin{equation}\label{eq:physical-final-edge}
 \mathcal A x_{t+L-1}=a_{t+L-1}x_{t+L}.
\end{equation}

\begin{corollary}[Best possible cyclic final column]\label{cor:final-column}
Among all cyclic final columns that land in the seed line $\mathbb Cx_t$,
the minimum physical residual is
\begin{equation}\label{eq:final-lower-bound}
 a_{t+L-1}d_{\mathbb P}(x_t,x_{t+L}).
\end{equation}
The best phase-only residual is
$a_{t+L-1}d_{\mathbb T}(x_t,x_{t+L})$.
\end{corollary}

Thus no alternative wrap normalization can beat
\eqref{eq:final-lower-bound} without changing the temporal span or the
physical final image.

The same argument identifies the minimal structural repair.  Let
$E\subset\cH$ be any proposed $k$-dimensional return space and let $P_E$ be
its orthogonal projection.

\begin{proposition}[Multi-edge return-space obstruction]
For a unit endpoint $v$ and amplitude $a\ge0$,
\begin{equation}\label{eq:return-subspace}
 \min_{z\in E}\norm{av-z}=a\norm{(I-P_E)v}.
\end{equation}
\end{proposition}

\begin{proof}
The orthogonal decomposition $av=P_E(av)+(I-P_E)(av)$ shows that the unique
best approximation is $z=P_E(av)$ and that the residual is the orthogonal
component in \eqref{eq:return-subspace}.
\end{proof}

The one-edge cyclic closure is the special case $E=\mathbb Cu$.  Therefore a
successful repair must enlarge or change the return space, alter the
temporal span, or use distinct right and left spaces.  Reweighting the same
seed line cannot help.

\section{Polar conditioning bounds}

Assume $G=J^*J>0$ and $V=JG^{-1/2}$.  Let $C(\omega)$ be the phase-marked
cycle and $K=G^{1/2}C(\omega)G^{-1/2}$.  The polar residual is
\begin{equation}\label{eq:polar-residual}
 R=\mathcal AV-VK
 =a_{t+L-1}(x_{t+L}-\omega x_t)e_{L-1}^*G^{-1/2}.
\end{equation}

\begin{proposition}[Two-sided Gram bound]\label{prop:gram-bound}
Let $\lambda_{\min}$ and $\lambda_{\max}$ be the extreme eigenvalues of
$G$.  Then
\begin{equation}\label{eq:gram-bound}
 \frac{a_{t+L-1}d_{\mathbb T}}
 {\sqrt{\lambda_{\max}}}
 \le \norm R\le
 \frac{a_{t+L-1}d_{\mathbb T}}
 {\sqrt{\lambda_{\min}}}.
\end{equation}
\end{proposition}

\begin{proof}
Equation \eqref{eq:polar-residual} is a rank-one operator.  Its norm equals
$a_{t+L-1}d_{\mathbb T}\norm{e_{L-1}^*G^{-1/2}}$.  The norm of
$G^{-1/2}$ lies between $\lambda_{\max}^{-1/2}$ and
$\lambda_{\min}^{-1/2}$ on every unit vector, giving
\eqref{eq:gram-bound}.
\end{proof}

The upper bound explains why an ill-conditioned temporal Gram can amplify a
moderate endpoint error.  The lower bound shows that polar normalization
cannot make the defect disappear.

\section{Formula audit}

The deterministic audit uses dimensions $4,5,8,11$, lengths $3,4$, and ten
random trials per pair, for 80 cases.  It compares:
\begin{enumerate}[leftmargin=2em]
\item the direct matrix norm of the rank-one wrap with its closed formula;
\item the optimized arbitrary-scalar residual with
$a d_{\mathbb P}$;
\item the actual polar residual with both bounds in
\eqref{eq:gram-bound}.
\end{enumerate}
All 80 cases pass.  The maximum direct rank-one formula error is
$4.45\times10^{-16}$, the scalar formula error is zero at stored precision,
and every polar residual lies in its declared interval.

These tests are not physical evidence.  They certify the implementation of
the exact finite theorem.

\section{Replay of the physical windows}

The RH-182 archive contains 126 predeclared windows.  Forty-three have a
negative real endpoint correlation, so the orientation mark changes the
sign of the wrap.  Exactly those 43 windows show a nontrivial phase-mark
improvement.  The largest chord reduction is $1.5849$.

Despite that substantial correction,
\begin{equation}\label{eq:no-projective-return}
 \#\{d_{\mathbb P}\le0.25\}=0.
\end{equation}

The minimum is $0.37202$, leaving a finite gap $0.12202$ above the declared
threshold even after arbitrary scalar optimization.  Since
\eqref{eq:return-subspace} is an exact lower bound, Gram conditioning and
polar normalization can only redistribute or amplify the remaining defect;
they cannot erase its component orthogonal to the seed line.

\section{Stability of the obstruction under endpoint perturbations}

For unit pairs $(u,v)$ and $(\widetilde u,\widetilde v)$, Cauchy--Schwarz
gives
\begin{equation}\label{eq:correlation-stability}
 \left|
 |\langle u,v\rangle|-|\langle\widetilde u,\widetilde v\rangle|
 \right|
 \le \norm{u-\widetilde u}+\norm{v-\widetilde v}.
\end{equation}
Thus an outward enclosure of the two endpoint vectors immediately yields an
outward enclosure of the correlation and hence of both projective
distances.  This observation is useful for a later interval implementation:
one need not validate a nonlinear phase optimizer separately.  It suffices
to validate the endpoint balls and apply the closed formulas.

The present replay is not such an interval proof, so the finite margin is
reported as floating data.  The point of \eqref{eq:correlation-stability} is
methodological: the exact obstruction is compatible with rigorous outward
rounding should the same branch ever need a certified finite rejection.
The minimum projective distance is $0.3720$.  Hence even an arbitrary complex
final-edge scalar cannot meet the declared return threshold in any audited
window.  The orthogonal wrap/primal/adjoint success count remains zero.

The length-four windows at $\sigma=0.01$ illustrate the distinction.  Their
endpoint correlations are commonly negative and large in magnitude, so the
phase mark repairs a sign.  Yet a correlation magnitude near $0.9$ still
corresponds to a projective distance around $0.4$, and the adjoint residual
is larger still.  The issue is not a missing minus sign.

\section{What is rejected and what remains}

Within the fixed temporal span, the following rescue attempts are now
closed:
\begin{equation}\label{eq:closed-repairs}
 \text{positive wrap}
 \longrightarrow
 \text{phase wrap}
 \longrightarrow
 \text{arbitrary scalar wrap}.
\end{equation}
None reaches the physical gate.  A different start time was already included
in the 126-window audit.

The theorem does not constrain a left temporal space distinct from the right
one.  In a nonnormal problem, a Petrov--Galerkin pair can have small right and
left residuals without any single orthogonal reducing range.  The physical
source and observation provide canonical seeds for such a pair.  The next
paper therefore replaces one orthogonal frame by balanced biorthogonal
frames.

\section{Boundary}

This paper proves exact phase, scalar, projective, and Gram-conditioned wrap
formulas and reports their finite physical replay.  It does not prove an
all-level no-go theorem, reject lagged or biorthogonal clocks, validate a
complement resolvent, construct a physical Riesz shell, close interface R or
Gate A, or establish any Hilbert--Polya or Riemann-hypothesis statement.

\bibliographystyle{plainnat}
\bibliography{references}
\end{document}
